% Related Work / Literature Review
% Aligned with: MA thesis proposal (low-latency ML tracking, Architect game, single-leg balance).
% Audience: mixed (medical technology); cite from references/bibliography.bib only.
%
% EXPANSION GUIDE (target ~12--18 pages for ~80-page thesis):
% - Sec.2: Add "Scope of VR in rehab" (stroke, upper/lower limb); "Visual feedback modalities" (mirror vs avatar).
% - Sec.3: Add "Pose paradigms" (bottom-up vs top-down); "Filtering/smoothing" (jitter, 1-Euro); "Latency composition" (capture+inference+render); optional table RTMPose variants.
% - Sec.4: Add "Knee rehab context" (TKA/ACL); "Mirror therapy limitations"; "Engagement and compliance".
% - Sec.5: Add "Why action cameras fail"; "Pose requirements" (FPS, resolution, FOV).
% - Add 1--2 recent papers (2024--25) when you find them; extend Summary with one para on thesis structure if needed.

\chapter{Related Work}\label{ch:related}

This chapter reviews the technical and clinical foundations for a low-latency, gamified VR system for knee rehabilitation. The thesis aims to develop a motion tracking pipeline that drives an in-VR avatar with minimal delay (target end-to-end latency under 10\,ms for the tracking loop) and to evaluate a goal-oriented ``Architect'' game during single-leg balance exercises compared to standard visual feedback. The following sections address: (1) the role of VR in rehabilitation and the problem of simulator sickness; (2) real-time pose estimation and latency requirements; (3) gamification and motor learning principles; (4) sensing hardware and camera choice for the lab; and (5) a summary with the research gap.

% ---------------------------------------------------------------------------
\section{Virtual Reality in Rehabilitation}
% ---------------------------------------------------------------------------

\subsection{Scope and applications}

Virtual reality has moved from experimental use to established applications in physical rehabilitation. VR can provide controlled, repeatable environments for motor retraining, visual feedback of movement, and exposure to tasks that are difficult or unsafe to replicate in the clinic alone. In lower-limb and knee rehabilitation, systems range from high-end installations such as the Computer Assisted Rehabilitation Environment (CAREN), which uses motion capture and large screens to display a gender-neutral avatar, to consumer headsets with game-like exercises \cite{klausing2025feasibility}. Kim et al.\ \cite{kim2020vrrt} showed that virtual reflection therapy (viewing a virtual representation of the affected limb) improved balance and gait in chronic stroke patients more than traditional mirror therapy, supporting the use of avatar-based feedback for motor learning.

Applications span upper-limb retraining, stroke rehabilitation (balance and gait), and lower-limb and knee rehabilitation (e.g.\ after total knee arthroplasty or ACL reconstruction). Common benefits include repeatable task difficulty, objective feedback, and the possibility of home- or clinic-based use with consumer hardware. The present thesis focuses on knee rehabilitation and single-leg balance, where VR can deliver visual feedback of the lower limb and trunk to support proprioceptive training.

\subsection{Visual feedback modalities}

Feedback can be provided via a mirror (mirror therapy), a screen-based or head-mounted avatar driven by motion capture or pose estimation, or large-screen systems such as CAREN. Mirror therapy is limited by the need for an asymmetrical posture to view the reflection and by the fixed viewpoint. Avatar-based feedback avoids these constraints and can present a gender-neutral, third-person or first-person view of the body. Klausing et al.\ \cite{klausing2025feasibility} used a CAREN-style setup with 3D motion capture and a 180-degree screen to correct lower-extremity biomechanics (e.g.\ knee-dominant squat and dynamic Q-angle), with a majority of patients showing improved joint angles during squats. Such high-end systems are not always available; the thesis therefore targets a setup that uses a single RGB camera and ML-based pose estimation to drive an avatar in a head-mounted display.

\subsection{Simulator sickness and tolerability}

A major barrier to wider adoption is user discomfort. Simulator sickness (SS; symptoms such as nausea, disorientation, and oculomotor strain) is one of the main factors limiting duration and acceptance of VR therapy. Artificial movement (e.g.\ floating or teleportation in VR) and high end-to-end latency are known contributors. Hollaus et al.\ \cite{Hollaus.2024} developed a software tool to evaluate the tolerability of different VR movement types with minimal confounding from game-like content; tested with 111 participants, the tool showed no signs of inducing SS through content alone, providing a basis for assessing movement-induced discomfort and making VR-based treatments more accessible. Controlling latency and refresh rate (see Section~\ref{sec:latency}) is therefore critical both for presence and for reducing SS in rehabilitation applications.

% ---------------------------------------------------------------------------
\section{Motion Tracking and Latency}
\label{sec:latency}
% ---------------------------------------------------------------------------

Real-time pose estimation is the technical enabler for low-latency motion feedback in VR: the pipeline must turn raw camera frames into a stable skeleton representation fast enough that the user perceives the avatar as responding immediately to their movement. End-to-end motion-to-photon latency is the sum of (1) camera capture and transfer, (2) pose inference (and any filtering), and (3) rendering and display. A target of under 10\,ms for the tracking loop (capture + inference, possibly including light smoothing) leaves budget for rendering and display while staying below the 63\,ms cybersickness threshold \cite{stauffert2021latency}.

Capture latency can be reduced by minimising the number of buffered frames at the sensor or driver: when the application reads the ``newest'' frame instead of draining a multi-frame buffer, the delay between a real-world motion and the frame used for inference is smaller. On the inference side, latency can be traded against accuracy by (a) running the person detector only every $N$ frames and tracking pose in between (top-down pipelines), and (b) using smaller model variants (e.g.\ RTMPose-s or lightweight configurations) that sacrifice some accuracy for lower per-frame inference time. These design levers are relevant when targeting a sub-10\,ms tracking loop with consumer hardware; the methodology and implementation chapters describe how they are applied in the present system.

\subsection{Pose estimation frameworks}

Current systems for markerless pose estimation fall into top-down and bottom-up paradigms. Bottom-up methods first detect body parts or keypoints across the~image and then group them into person instances; they can be efficient in multi-person scenes but may sacrifice some accuracy. Top-down methods first detect the person (e.g.\ with a bounding box) and then estimate pose within that region. Top-down methods first detect the person (e.g.\ with a bounding box) and then estimate pose within that region, often achieving higher accuracy at the cost of more computation. The MMPose ecosystem and RTMPose models follow this approach. Jiang et al.\ \cite{Jiang.2023} present RTMPose, a real-time multi-person pose estimation framework that reaches 90+ FPS on an Intel i7-11700 CPU and 430+ FPS on a consumer GPU (NVIDIA GTX 1660 Ti), with RTMPose-m achieving 75.8\% AP on COCO. On high-end GPUs (e.g.\ RTX 4090), end-to-end latency is reported in the 7--10\,ms range; on mobile (Snapdragon 865), RTMPose-s still achieves 70+ FPS (~14\,ms per frame), which is relevant for keeping motion-to-display latency low in VR applications. Lugaresi et al.\ \cite{Lugaresi.2019} introduce MediaPipe, a framework for building perception pipelines including BlazePose, with efficient deployment across platforms and a detector-tracker design that reduces per-frame cost once the person is localised. For the present thesis, MMPose/RTMPose is used as the primary inference engine to meet the proposed latency target. Table~\ref{tab:rtmpose-variants} summarises representative RTMPose variants and their reported latency and throughput on different hardware; the 7--14\,ms range for inference (or inference plus detection) is compatible with a sub-10\,ms tracking loop when capture latency is low.

\begin{table}[htbp]
\centering
\caption{Representative RTMPose variants: latency and throughput (from \cite{Jiang.2023} and literature).}
\label{tab:rtmpose-variants}
\begin{tabular*}{\textwidth}{@{\extracolsep{\fill}}lllll@{}}
\toprule
Variant & Backbone & Hardware & Latency (ms) & Throughput (FPS) \\
\midrule
RTMPose-m & CSPNeXt-m & RTX 4090 & 7--10 & $\sim$100--140 \\
RTMPose-m & CSPNeXt-m & Intel i7-11700 & $\sim$11 & $\sim$90 \\
RTMPose-s & CSPNeXt-s & Snapdragon 865 & $\sim$14 & 70+ \\
\bottomrule
\end{tabular*}
\end{table}

\subsection{Filtering and smoothing}

Raw pose outputs from ML models can exhibit high-frequency jitter, which is undesirable for clinical range-of-motion assessment and for a stable avatar. MediaPipe-based pipelines in particular are often reported to need smoothing. Common approaches include a 10th-order Butterworth low-pass filter or the 1-Euro filter; the latter adapts to movement speed to minimise lag during fast motion while keeping slow, static poses smooth. For the Architect scenario (single-leg balance), where movements are relatively slow, light filtering can be applied without significantly increasing effective latency, as long as the filter is tuned to preserve responsiveness.

\subsection{Latency and cybersickness}

In VR, the delay between the user's motion and the update of the visual display (motion-to-photon latency) directly affects presence, motor performance, and comfort. Stauffert et al.\ \cite{stauffert2021latency} showed that end-to-end latency above 63\,ms induces significant cybersickness symptoms in immersive VR; motor performance and the perception of simultaneity degrade by about 75\,ms. Waltemate et al.\ \cite{waltemate2016impact} demonstrated in a CAVE setup with full-body avatars that the sense of body ownership and agency breaks down when tracking latency exceeds approximately 125\,ms. Table~\ref{tab:latency-thresholds} summarises thresholds reported in the literature; these motivate the thesis goal of a tracking loop under 10\,ms so that total motion-to-photon latency remains well below the 63\,ms comfort limit.

\begin{table}[htbp]
\centering
\caption{Reported latency thresholds relevant to VR rehabilitation (motion-to-photon or tracking delay).}
\label{tab:latency-thresholds}
\begin{tabularx}{\textwidth}{@{}lXl@{}}
\toprule
Approx.\ threshold & Effect & Source \\
\midrule
16--20\,ms & Perceptual threshold (delay noticeable in fast scenarios) & \cite{ellis2004generalizeability} \\
63\,ms & Significant cybersickness (SSQ) & \cite{stauffert2021latency} \\
75\,ms & Motor performance \& simultaneity impaired & \cite{stauffert2021latency} \\
125\,ms & Body ownership \& agency begin to break down & \cite{waltemate2016impact} \\
\bottomrule
\end{tabularx}
\end{table}

Frame rate also affects comfort: higher frame rate reduces frame time and perceptual mismatch. Industry and usability reports suggest that 72\,Hz (or higher) significantly improves tolerance compared to 60\,Hz in active VR scenarios \cite{battlestart2024fps}. For rehabilitation, targeting 72--90\,Hz or higher is advisable so that patients can focus on the exercise rather than discomfort.

% ---------------------------------------------------------------------------
\section{Gamification and Motor Learning}
% ---------------------------------------------------------------------------

The ``Architect'' game is designed to support single-leg balance and proprioceptive training (e.g.\ after total knee arthroplasty or ACL reconstruction) by shifting the patient's focus from the joint itself to an external goal (e.g.\ keeping a virtual structure stable). This aligns with the external focus of attention (EFA) principle in motor learning.

\subsection{Knee rehabilitation context}

After total knee arthroplasty (TKA) or ACL reconstruction, patients often experience reduced proprioception and balance due to surgical disruption of ligamentous and capsular mechanoreceptors. Rehabilitation typically combines strength training with balance and coordination tasks (e.g.\ single-leg stance, tandem walking, figure-of-eight walks). Evidence suggests that proprioception-focused programmes lead to better functional outcomes than strength-only programmes, supporting the use of balance-oriented VR tasks such as single-leg standing in the Architect game.

\subsection{Proprioception and balance after knee surgery}

Proprioception (afferent information from mechanoreceptors integrated for postural control) is often impaired after TKA due to disruption of ligamentous and capsular structures. Jogi et al.\ \cite{jogi2024proprioception} conducted a prospective randomised study comparing proprioception-based rehabilitation (balance and coordination tasks) with routine strength-focused exercise in TKA patients. At six months, the proprioception group showed significantly better outcomes on the 6-minute walk, chair stand, fast walk, and Timed Up and Go tests, supporting the emphasis on balance and coordination in the Architect scenario rather than on strength alone.

\subsection{External focus of attention}

The constrained action hypothesis suggests that an internal focus (attention on one's own body) increases conscious control and can interfere with automatic motor patterns, whereas an external focus (attention on the effect of the movement in the environment) promotes more efficient, stable performance. Sherman et al.\ \cite{sherman2021external} used EEG during single-leg balance under internal vs.\ external focus instructions and found that EFA was associated with cortical activity consistent with lower cognitive effort and greater automaticity. Gamifying balance by having the user stabilise a virtual object (the ``Architect'' mechanic) naturally induces an external focus and may improve retention and transfer to real-world balance tasks \cite{jogi2024proprioception}.

\subsection{Mirror therapy and VR-based alternatives}

Traditional mirror therapy uses a mirror to create the illusion that the affected limb is moving correctly. Limitations include the need for an asymmetrical viewing posture (which can interfere with balance training) and a fixed viewpoint. Kim et al.\ \cite{kim2020vrrt} showed that virtual reality reflection therapy (VRRT; viewing a virtual representation of the limb via a camcorder and screen) improved postural balance and gait in chronic stroke patients more than mirror therapy, and was measured with Berg Balance Scale and Timed Up and Go. This supports the thesis approach of using an avatar driven by pose estimation instead of a mirror for single-leg balance feedback.

\subsection{Engagement and compliance}

Conventional knee exercises often suffer from low patient engagement and adherence. Goal-oriented, game-like tasks can increase motivation and time-on-task. The Architect mechanic (e.g.\ keeping a virtual structure stable while balancing) provides a clear external goal and implicit feedback, which may support both motor learning and adherence. The pilot evaluation in this thesis will assess engagement and perceived responsiveness (system feel) compared to standard visual feedback.

% ---------------------------------------------------------------------------
\section{Sensing Hardware and Camera Choice}
% ---------------------------------------------------------------------------

\subsection{Why action cameras are unsuitable for low-latency VR}

Action cameras such as the GoPro Hero 4 are designed for high-quality recording rather than real-time streaming. Over USB or Wi-Fi, reported glass-to-glass latency (sensor to displayed image) can reach 210\,ms or more, and in some configurations 500\,ms. For a VR loop targeting under 30\,ms motion-to-photon, such delays are unacceptable: they exceed the 63\,ms cybersickness threshold and the 125\,ms threshold at which body ownership of the avatar breaks down \cite{waltemate2016impact}. HDMI output from the camera can reduce latency to roughly 80--150\,ms but is often unstable or impractical for integration with a PC-based pose pipeline. Therefore, the thesis and lab setup rely on standard USB video devices rather than action cameras.

\subsection{Requirements for pose estimation and lab setup}

Pose estimation models typically resize input to modest resolutions (e.g.\ 256$\times$256 or 384$\times$384); thus very high sensor resolution is less important than high frame rate and low capture latency. A camera that delivers 720p or 1080p at 60\,FPS over USB 3.0 (UVC) usually adds only a few milliseconds to the pipeline. Field of view should cover the user's full body (or at least lower body and trunk) when standing in front of the sensor. For the Architect scenario, a single frontal or slightly angled camera is sufficient; multi-camera setups can improve robustness but increase cost and complexity.

\subsection{Sensor options and recommendation}

Action cameras (e.g.\ GoPro Hero 4) are optimised for recording rather than live streaming; reported glass-to-glass latency over USB or Wi-Fi can reach 200\,ms or more, which is unsuitable for VR feedback. In contrast, standard USB 3.0 webcams can deliver 720p or 1080p at 60\,FPS with much lower capture latency (often under 30--50\,ms for the camera subsystem). Martínez-García et al.\ \cite{mdpi2022exercam} demonstrated that a simple RGB webcam with real-time pose estimation (OpenPose and similar) is feasible for telerehabilitation and acceptable in home environments, supporting the use of consumer-grade cameras in the lab. Table~\ref{tab:camera-options} summarises sensor options relevant to the thesis and to lab procurement (e.g.\ budget under 200\,EUR for additional hardware if needed).

\begin{table}[htbp]
\centering
\caption{Sensor options for real-time pose estimation in VR rehab (lab use; approximate prices).}
\label{tab:camera-options}
\begin{tabularx}{\textwidth}{@{}lXlX@{}}
\toprule
Sensor / type & Interface & Typical latency & Suitability \\
\midrule
GoPro Hero 4 (USB/Wi-Fi) & NCM, MPEG-TS & 210--500\,ms & Not suitable for VR \\
Standard USB 3.0 webcam (720p60) & UVC & $<$30\,ms & Suitable \\
Logitech Brio / C920-class & UVC & $<$50\,ms & Suitable \\
OAK-1 Lite (on-device AI) & USB 3 & Low (on-device) & Optional ($\sim$99) \\
\bottomrule
\end{tabularx}
\end{table}

\noindent\textit{Abbreviations:} UVC = USB Video Class; NCM = Network Control Model.

For the present project, a standard RGB camera (e.g.\ USB 3.0, 60\,FPS or higher) is the primary recommendation; if the lab considers a dedicated vision module, the OAK-1 Lite offers on-device inference at low latency within the stated budget. Existing lab equipment that meets UVC 60\,FPS and low-latency capture is compatible with the MMPose/RTMPose pipeline.

% ---------------------------------------------------------------------------
\section{Summary and Research Gap}
% ---------------------------------------------------------------------------

Existing VR rehabilitation setups often rely on tracking solutions with high end-to-end latency (on the order of 200\,ms), leading to laggy avatar behaviour, reduced presence, and increased risk of simulator sickness. At the same time, conventional knee rehabilitation exercises frequently suffer from low engagement and compliance. This thesis addresses both issues by (1) developing a low-latency ML-based motion tracking system using a standard RGB camera and MMPose/RTMPose, with a target tracking loop under 10\,ms, and (2) implementing a gamified ``Architect'' prototype for single-leg balance that uses an external focus of attention.

Prior work covers real-time pose estimation \cite{Jiang.2023, Lugaresi.2019}, latency and cybersickness thresholds \cite{stauffert2021latency, waltemate2016impact}, VR-based reflection and feedback for rehabilitation \cite{kim2020vrrt, klausing2025feasibility}, proprioception-focused knee rehab \cite{jogi2024proprioception}, and the role of external focus in balance \cite{sherman2021external}. Tools to evaluate VR movement tolerability and simulator sickness have been developed in the same research context \cite{Hollaus.2024}. However, there is no reported system that combines \emph{low-latency} (sub-30\,ms, targeting sub-10\,ms for the tracking component) \emph{markerless} skeletal tracking via a consumer RGB camera with a \emph{gamified}, goal-oriented single-leg balance task in VR for knee rehabilitation, and that is evaluated for both system feel (responsiveness) and engagement compared to standard visual feedback. This gap is what the present thesis aims to fill through the development of the tracking pipeline, the Architect prototype, and a pilot user evaluation. The following chapters describe the methodology, implementation, and evaluation that address this gap.
